
\section{Datasets}
\label{sec-datasets}
This section introduces the datasets, including their background TKG, size, etc. We provide a question category coverage comparison across TKGQA datasets in Table~\ref{questincategoryfordataset}.


\subsection{TempQuestions} TempQuestions \cite{jiaTempQuestionsBenchmarkTemporal2018} is a benchmark dataset derived from Freebase~\cite{bollackerFreebaseCollaborativelyCreated2008}, where temporal knowledge is stored using compound value types (CVTs). Examples of these CVTs include \texttt{footballPlayer.team.joinedOnDate}, \texttt{footballPlayer.team.leftOnDate},
\texttt{marriage.date},
\texttt{amusement\_parks.ride.opened}	,\texttt{amusement\_parks.ride.closed}.
It includes 1,271 questions with temporal signals, question types, and data sources for testing and evaluation.

\subsection{\textbf{TempQA-WD}} TempQA-WD~\cite{neelamSYGMASystemGeneralizable2021} 
is an adaptation of TempQuestions for Wikidata~\cite{vrandecicWikidataFreeCollaborative2014}, fulfilling the need for a multi-KG-based dataset to evaluate their model.
The dataset comprises:
\begin{itemize}
    \item 839 questions with corresponding Wikidata SPARQL queries, answers, categories, and TempQuestions' information.
    \item 175 questions with AMR, lambda expressions, entities, relations, and KB-specific lambda expressions, in addition to the above information.
\end{itemize}


\subsection{\textbf{TimeQuestions}}TimeQuestions~\cite{jiaComplexTemporalQuestion2021}
is based on Wikidata, it includes temporal facts of triples, such as 
\texttt{<Malia Obama, date of birth, 04-07-1998>} or maintain more temporal knowledge with qualifiers like
\texttt{<Barack Obama, position held, President of the US; start date, 20-01-2009; end date, 20-01-2017>}. 
Jia searched through eight KG-QA datasets for time-related questions and mapped them to Wikidata. Questions in each benchmark are tagged for temporal expressions using SUTime~\cite{changSUTIMELibraryRecognizing} and HeidelTime~\cite{strotgenHeidelTimeHighQuality2010} and for signal words using a dictionary compiled by Setzer~\cite{setzerTemporalInformationNewswire2001a} and manually tagged with its temporal question category.
In total, the TimeQuestions comprises 16,859 questions.

\subsection{\textbf{CronQuestions}}CronQuestions~\cite{chenTemporalKnowledgeGraph2022} utilises a subset of Wikidata that includes facts annotated with temporal information~\cite{lacroixTensorDecompositionsTemporal2020}, such as \texttt{<Barack Obama, held position, President of USA, 2008, 2016>}. Entities extracted from Wikidata with both ``start time'' and ``end time'' annotations are transformed into event format (e.g., \texttt{<WWII, significant event, occurred, 1939, 1945>}). 
The dataset comprises a Temporal KG with 125k entities and 328k facts (including 5k event facts), and 410k natural language questions requiring temporal reasoning.
% (subject， relation， object， timestamp)
% "To generate the QA dataset, we started with templates for temporal reasoning."
%"These were made using the five most frequent relations from our WikiData subset, namely • member of sports team • position held • award received • spouse • employer"
%"We then asked human annotators to paraphrase these templates to generate more linguistic diversity."


\subsection{\textbf{Complex-CronQuestions}}
Chen~\cite{chenTemporalKnowledgeGraph2022} observe that existing benchmarks contain many pseudo-temporal questions. For instance, for the question \textit{``What’s the first award Carlo Taverna got?''} There is only one fact related to\textit{ Carlo Taverna} in the TKG, which makes the temporal word ``first'' meaningless as a constraint. 
They remove all simple and pseudo-temporal questions and filter out questions with fewer than 5 relevant facts in CronQuestions.


\subsection{\textbf{MultiTQ}}MultiTQ \cite{chenMultigranularityTemporalQuestion2023} is a dataset derived from ICEWS05-15~\cite{garcia-duranLearningSequenceEncoders2018}, where all facts are standardized as quadruple $(s,r,o,t)$.
ICEWS05-15 is notable for its rich semantic information, with a higher average number of relation types per entity than other TKGs.
The MultiTQ dataset features several advantages, including its large scale, ample relations, and multiple temporal granularities.
ICEWS provides time information at a day granularity, while the authors generate higher granularities, such as year and month, for the questions. MultiTQ contains 500,000 questions, making it a significant resource for temporal question-answering research.

\subsection{\textbf{MusTQ}}MusTQ \cite{zhangMusTQTemporalKnowledge2024} is a dataset that contains diverse multi-step temporal reasoning. MusTQ contains ~666K temporal questions and a WikiData subset (with 125k entities and 326k facts) as its knowledge base. It has six question types and four answer types (entity, time, boolean, numeric), and the maximum number of reasoning steps required for the questions is up to four.


\subsection{\textbf{ForecastTKGQuestions}}ForecastTKGQuestions \cite{dingForecastTKGQuestionsBenchmarkTemporal2022} aims to predict facts in the future timestamps beyond the ICEWS subset. It includes three types of forecasting questions, i.e., entity prediction, yes-unknown, and fact reasoning questions. The answer types for each category are entity, yes-unknown label, and choice label, respectively. The total number of questions in ForecastTKGQuestions is up to 727k. Entities and timestamps are annotated in their questions. 

\subsection{\textbf{TIQ}}TIQ~\cite{jiaFaithfulTemporalQuestion2024} is a benchmark with a primary focus on challenging and diverse implicit questions, and it has multiple sources: Wikipedia text and infoboxes, and the Wikidata KB. It has 10,000 questions with metadata. The metadata includes the information snippets grounding the question, the sources these were obtained, the explicit temporal value expressed by the implicit constraint, the topic entity, the question entities detected in the snippets, and the temporal signal.


% \begin{figure*}[h]
%     \centering
%     \includegraphics[width=0.9\textwidth]{Category.pdf}
%     % 已有工作中涉及到的问题分类
%     \caption{Temporal Question category with examples.}
%     \label{fig:pdfimage}
%     % 缩一缩，把附录的表缩一下，往上放。
%     % 把这个放附录。
% \end{figure*}

\begin{table*}[htbp]
    \caption{Question category coverage comparison across TKGQA datasets.\label{questincategoryfordataset}}
    \centering
    \scriptsize
    \renewcommand{\arraystretch}{1.8}
    \resizebox{\linewidth}{!}{
    \begin{NiceTabular}{|c|c|c|c|c|c|c|c|c|c|}
    \toprule
        \Block{1-3}{\textbf{Category}}& & & \multicolumn{1}{m{1.1cm}}{\textbf{\shortstack{Temp\\Questions}}} & \multicolumn{1}{m{1.1cm}}{\textbf{\shortstack{TempQA-\\WD}}} & \multicolumn{1}{m{1.1cm}}{\textbf{\shortstack{Time\\Questions}}} & \multicolumn{1}{m{1.1cm}}{\textbf{\shortstack{Cron\\Questions}}} & \multicolumn{1}{m{1.6cm}}{\textbf{\shortstack{Complex\\CronQuestions}}} & \textbf{MultiQA} & \textbf{MusTQ}\\ \hline
        \multirow{13}{*}{Question Content} & \multirow{3}{*}{\shortstack{Time\\Granularity}} & Year & \checkmark & \checkmark & \checkmark & \checkmark &\checkmark& \checkmark &\\ 
        ~ & ~ & Month & \checkmark & \checkmark & ~ & ~ & & \checkmark & \checkmark \\
        ~ & ~ & Day & \checkmark & \checkmark & \checkmark & ~ && \checkmark & \\  \cline{2-10}
        ~ & \multirow{2}{*}{\shortstack{Time\\Expression}} & Explicit & \checkmark & \checkmark & \checkmark & \checkmark &\checkmark& \checkmark & \\ 
        ~ & ~ & Implicit & \checkmark & \checkmark & \checkmark & \checkmark & \checkmark&\checkmark & \checkmark \\ \cline{2-10}
        ~ & \multirow{6}{*}{\shortstack{Temporal\\Signal}} & Overlap & \checkmark & \checkmark & \checkmark & \checkmark &\checkmark &\checkmark & \checkmark \\ 
        ~ & ~ & Equal & \checkmark & \checkmark & \checkmark & \checkmark &\checkmark &\checkmark & \checkmark \\ 
        ~ & ~ & Start/End & \checkmark & \checkmark & \checkmark & \checkmark &\checkmark & ~ & \\ 
        ~ & ~ & During/Include & \checkmark & \checkmark & \checkmark & \checkmark &\checkmark & \checkmark & \checkmark\\ 
        ~ & ~ & Before/After & \checkmark & \checkmark & \checkmark & \checkmark &\checkmark &\checkmark & \checkmark\\ 
        ~ & ~ & Ordinal & \checkmark & \checkmark & \checkmark & \checkmark & \checkmark&\checkmark & \checkmark\\ \cline{2-10}
        ~ & {\multirow{2}{*}{\shortstack{Temporal\\Constraint\\Composition}}} & w/ Composition & & ~ & ~ & ~ & &\checkmark & \checkmark\\
        ~ & ~ & w/o Composition & \checkmark & \checkmark & \checkmark & \checkmark & \checkmark &\checkmark & \checkmark\\ \hline
      \multirow{5}{*}{Answer Type} & \Block{1-2}{Entity} & ~  & \checkmark & \checkmark & \checkmark & \checkmark &\checkmark &\checkmark & \checkmark \\ \cline{2-10} 
        ~ & \multirow{3}{*}{\shortstack{Time}} & Year & \checkmark & \checkmark & ~ & \checkmark &\checkmark &\checkmark & \checkmark \\ 
        ~ & ~ & Month & \checkmark & \checkmark & ~ & ~ & &\checkmark & \\ 
        ~ & ~ & Day & \checkmark & \checkmark & \checkmark & ~ & &\checkmark & \\  \cline{2-10}
         & \Block{1-2}{Boolean} & & ~ & ~ & ~ & ~ & & ~ & \checkmark \\  \hline
        \multirow{2}{*}{Complexity} & \Block{1-2}{Simple} & &  \checkmark & \checkmark & \checkmark & \checkmark & &\checkmark & \checkmark\\ \cline{2-10}
        ~ &  \Block{1-2}{Complex} & & \checkmark &  \checkmark  &\checkmark & \checkmark & \checkmark &\checkmark & \checkmark\\ 
        \bottomrule
    \end{NiceTabular}}
\end{table*}

\section{Decomposition-based TKG Embedding Methods}
\label{TKGembedding}
Temporal Knowledge Graph Embedding attempts to learn the representations of entities, relations, and timestamps in low-dimensional latent feature spaces while preserving certain properties of the original graph.

Formally speaking, given a TKG $\mathcal{K} = (\mathcal{E,R,T,F})$, TKG embedding methods typically learn $D$-dimential vectors $\mathbf{e}_\epsilon, \mathbf{v}_r, \mathbf{t}_\tau \in \mathbb{R}^D$ for each $\epsilon \in \mathcal{E}$, $r\in\mathcal{R}$ and $\tau\in \mathcal{T}$. These emebdding vectors are learned such that each valid fact $(s,r,o,\tau) \in \mathcal{F}$ is scored higher than an invalid fact $\left(s^{\prime}, r^{\prime}, o^{\prime}, \tau^{\prime}\right) \notin \mathcal{F}$ through a scoring function $\phi(·)$, i.e., $\phi\left(\mathbf{e}_s, \mathbf{v}_r, \mathbf{e}_o, \mathbf{t}_\tau\right)>\phi\left(\mathbf{e}_{s^{\prime}}, \mathbf{v}_{r^{\prime}}, \mathbf{e}_{o^{\prime}}, \mathbf{t}_{\tau^{\prime}}\right)$~\cite{kazemiRepresentationLearningDynamic2020}. Different TKG embedding models have been proposed so far to efficiently learn the representations of TKGs and perform TKG completion as well as inference~\cite {zhangSurveyTemporalKnowledge2024b}.

Decomposition-based methods are one line of the TKG embedding methods. A tensor is a multidimensional array~\cite{kolda2009tensor}.
Tensor decomposition has three applications: dimension reduction, missing data completion, and implicit relation mining, which meet the needs of representation learning.~\cite{caiSurveyTemporalKnowledge2024}

For the temporal knowledge graph, the quadruple can be represented by an order of four tensors, and each tensor dimension is the head entity, relation, tail entity, and timestamp, respectively.

\subsection{TComplex and TNTComplex}
Lacroix et al.~\cite{lacroixTensorDecompositionsTemporal2020} utilise tensor decomposition and proposed the TComplEx and TNTComplEx based on ComplEx~\cite{trouillonComplexEmbeddingsSimple2016}. 

ComplEx is a KGE method. For each fact $(s, r, o)$ in static KG, all entities and relations are embedded into a Complex vector space $\mathbb{C}$. In detail, ComplEx represents facts by three sets of embedding vectors: one for head entities, another for tail entities, and a third for relations. The score function $\phi (s, r, o)$ is defined as: 

$$
\begin{aligned}
\phi(s, r, o) & =\mathfrak{R}\left(<\boldsymbol{e}_s, \boldsymbol{v}_r, \overline{\boldsymbol{e}}_o>\right) \\
& =\mathfrak{R}\left(\sum_{i=1}^d \boldsymbol{e}_{s_i} \boldsymbol{v}_{r_i} \overline{\boldsymbol{e}}_{o_i}\right)
\end{aligned}
$$

where $\bar{\boldsymbol{e}}$ represents the complex conjugate of $\boldsymbol{e}$, $\mathfrak{R}$ denotes adopting the real part of the complex number in the score function.


TComplEx extended ComplEx by increasing the dimensions of the fact tensor. The timestamp of temporal facts is also denoted as complex vector $t_\tau \in \mathbb{C}^{d\times1}$. The score function  $\phi(s, r, o, \tau)$ is transformed as follows:

$$
\begin{array}{r}
\phi(s, r, o, \tau)=\mathfrak{R}\left(<\boldsymbol{e}_s, \boldsymbol{v}_r, \overline{\boldsymbol{e}}_o, \boldsymbol{t}_\tau>\right) \\
=\mathfrak{R}\left(<\boldsymbol{e}_s, \boldsymbol{v}_r \odot \boldsymbol{t}_\tau, \overline{\boldsymbol{e}}_o>\right)
\end{array}
$$
$\odot$ denotes the element-wise product.  

TComplEx employs additional regularizers to improve the quality of the learned embeddings, such as enforcing close timestamps to have similar embeddings (closeness of time). The embedding learning procedure makes TComplEx a suitable method for inferring missing facts such as $(s, r, ?, t)$ and $(s, r, o, ?)$ over an incomplete TKG. Due to its aforementioned benefits, it is used by most TKG Embedding-based TKGQA methods.


Lacroix et al.~\cite{lacroixTensorDecompositionsTemporal2020} also discovered that in the real world, some relationships between entities remain constant, such as ``is located in'', where Shanghai always remains in China. Therefore, TNTComplEx introduces the non-temporal part together with TComplEx. It decomposes a tensor into a sum of tensors with and without time information. This allows the sharing of parameters between parts with and without time information. To represent adjacent timestamps more closely, a penalty is applied to the norm of time embeddings using the smoothness theory.

The score function for TNTComplEx denotes as $\mathfrak{R}\left(<\boldsymbol{e}_s, \boldsymbol{v}_r, \overline{\boldsymbol{e}}_o, \boldsymbol{t}_\tau> + e_s, <\boldsymbol{e}_s, \boldsymbol{u}_r, \overline{\boldsymbol{e}}_o, 1>\right)$, where $ \boldsymbol{u}_r$ represents the temporal agnostic relation representation.

Many other effective TKGE methods may be applied to the TKGQA task~\cite{liuTLogicTemporalLogical2022a,yuanTRHyTETemporalKnowledge2022, xiaMetaTKGLearningEvolutionary2023}, leaving it for the reader to explore.

